\begin{abstract}\label{chp:abstract}
\vspace{-1.5em}

A major challenge in large-scale GPU clusters is that schedulers cannot predict or control which heterogeneous, heavy-tailed jobs are assigned to shared hardware. Because of the long-tailed nature of workloads, GPU clusters experience extreme head-of-line blocking when a large ``elephant'' blocks many smaller ``mouse'' tasks. The traditional Predict-then-Optimize (PtO) framework attempts to solve this by minimizing point prediction errors (e.g., mean absolute error) before scheduling. In this dissertation, we propose a decision-focused learning (DFL) approach to provide a new direction in machine learning. We show that ultimately, the objective of the predictor is not just to achieve accurate predictions on its own merit, but to optimize overall system-level performance measures, specifically Job Completion Time (JCT). Using data from Alibaba's Platform for Artificial Intelligence (PAI), we developed an end-to-end testbed tying runtime predictions to event-driven simulations across both small (32-GPU) and massive (256-GPU) cluster environments. Each of our six models is being used to predict performance as measured by a new ``cluster utilization'' method using a ``sweep-line''. Results of the analysis show that there is a very dramatic difference in how well the predictive capability of the algorithms works depending on the size of the cluster. Specifically, the tree-based algorithms (XGBoost and LightGBM) produced much lower MAEs, which resulted in significantly better optimization of JCT. Deep learning models (categorical LSTMs) produce higher MAEs; however, they act as superior ``Learning to Rank'' engines protecting the ``mice'' from the ``elephants.'' Also, the LSTM produces up to a 57.25\% reduction in JCT over the baseline. These results provide concrete evidence that at least in massive-scale cluster scheduling, making the right relative ranking decisions is far more important than achieving the right absolute prediction.

\end{abstract}